
\section{Uncertainty Quantification}
\label{sec:uq}

% ---------------------------------------------------------------
% Sub-topics to cover:
%
% 1. Why UQ matters in HEP
%    - Physics measurements require not just a central value but
%      a rigorous uncertainty.  ML methods add a new source of
%      uncertainty (model uncertainty) on top of statistical and
%      systematic uncertainties.
%    - The ATLAS top tagger paper (Ch.4) deals explicitly with
%      how to assign and propagate ML model uncertainties.
%
% 2. Types of uncertainty
%    - Aleatoric (irreducible, data noise) vs. epistemic (reducible,
%      model / training uncertainty).
%    - In HEP: statistical uncertainty (finite data) is aleatoric;
%      choice of MC generator, network architecture, training
%      procedure are epistemic.
%
% 3. Ensemble methods
%    - Train N models with different seeds / architectures.
%    - Variance of ensemble predictions estimates epistemic uncertainty.
%    - Used in the top tagger paper: ensemble of taggers to
%      derive model uncertainty on the WP efficiency.
%
%
% 5. Systematic uncertainties from ML
%    - Variation of training hyperparameters.
%    - Dependence on MC generator (Pythia vs. Herwig).
%    - "Generator uncertainty" is a major theme in Ch.4 and Ch.6.
